\section{Quality of Generated Code}

This section applies specification satisfaction --- whether a feature's behavior fulfills its functional and non-functional requirements and handles its edge cases, not merely whether its tests pass --- to Part 1 and then to Part 2.

Part 1's functional requirements were generally met, but rarely on the first attempt and rarely without a gap the prompt itself had left open. The missing email field, the missing uniqueness constraint on it, and the undifferentiated validation messages in the login and registration feature were each caught only once the implementation was reviewed against what the feature was actually meant to do, not against what had literally been asked for --- meaning specification satisfaction in Part 1 depended on a human noticing the gap after the fact, feature by feature, rather than on anything in the process that would have surfaced it beforehand. Non-functional requirements fared worse. The unauthenticated \textbf{\texttt{user\_id}} parameter was a direct failure of the security requirement --- any user could view or modify another user's data simply by changing a number in the URL. Unlike the login-feature gaps, which were caught and fixed once someone reviewed the output, this one was never caught during Part 1 at all. The three architectural weaknesses --- the missing service layer, the MainPage.js monolith, and the scattered authentication reads --- represent the same failure at the level of maintainability: each satisfied its immediate functional requirement while working against the maintainability the non-functional requirements called for, and none of the three was caught until its consequences had already spread across multiple files. Part 1's code, in short, tended to satisfy what a feature was explicitly asked to do and to fail, sometimes for the whole duration of the phase, at the requirements that were assumed rather than stated.

Part 2 shifted where that gap could still occur. The mandatory Design Review moved the check earlier --- the super-admin plan's privacy concern was identified and given a mitigation while the feature existed only as a document, not after code already existed to review --- and several Charter rules encode a specific Part 1 miss directly, so that the same non-functional failure could not recur silently. Because Part 2's added scope was structurally more varied than Part 1's --- access control, multi-user collaborative content, cross-subsystem logging, historical aggregation --- specification satisfaction had to hold across a wider range of requirement types, not just repeated instances of the same one. The inherited weaknesses, including the unresolved \textbf{\texttt{user\_id}} vulnerability, were not left as background debt: each went through the same Architect--human--Implementer--Reviewer sequence as new work, and a separate, independent security audit added a second check specifically aimed at the class of non-functional requirement Part 1 had failed on. Part 2's code, on the whole, satisfied both its functional and its non-functional requirements more consistently, and --- where the Design Review gate worked as intended --- pre-emptively, rather than only once a gap had already become visible in the running application.

Across both parts, specification satisfaction moved from something checked after code existed to something checked, at least for the highest-stakes cases, before it did. Part 1 satisfied what was explicitly asked and missed what was assumed, sometimes for the full duration of the phase. Part 2 caught more of what would have been assumed earlier, and treated inherited gaps as work to be done rather than debt to be carried forward.

